\documentclass[10pt,a4paper]{amsart}
\usepackage{amsmath,amssymb,amsthm,graphicx}
\usepackage{verbatim}
\newenvironment{reference}{\comment}{\endcomment}
%\newenvironment{reference}{}{}
\newenvironment{slogan}{\comment}{\endcomment}
\newenvironment{history}{\comment}{\endcomment}
% For commutative diagrams we use Xy-pic
\usepackage[all]{xy}

% We use 2cell for 2-commutative diagrams.
\xyoption{2cell}
% Theorem environments.
%
\theoremstyle{plain}
\newtheorem{theorem}[subsection]{Theorem}
\newtheorem{proposition}[subsection]{Proposition}
\newtheorem{lemma}[subsection]{Lemma}

\theoremstyle{definition}
\newtheorem{definition}[subsection]{Definition}
\newtheorem{example}[subsection]{Example}
\newtheorem{exercise}[subsection]{Exercise}
\newtheorem{situation}[subsection]{Situation}

\theoremstyle{remark}
\newtheorem{remark}[subsection]{Remark}
\newtheorem{remarks}[subsection]{Remarks}

\numberwithin{equation}{subsection}

% Macros
%
\def\lim{\mathop{\mathrm{lim}}\nolimits}
\def\colim{\mathop{\mathrm{colim}}\nolimits}
\def\Spec{\mathop{\mathrm{Spec}}}
\def\Hom{\mathop{\mathrm{Hom}}\nolimits}
\def\Ext{\mathop{\mathrm{Ext}}\nolimits}
\def\SheafHom{\mathop{\mathcal{H}\!\mathit{om}}\nolimits}
\def\SheafExt{\mathop{\mathcal{E}\!\mathit{xt}}\nolimits}
\def\Sch{\mathit{Sch}}
\def\Mor{\mathop{\mathrm{Mor}}\nolimits}
\def\Ob{\mathop{\mathrm{Ob}}\nolimits}
\def\Sh{\mathop{\mathit{Sh}}\nolimits}
\def\NL{\mathop{N\!L}\nolimits}
\def\CH{\mathop{\mathrm{CH}}\nolimits}
\def\proetale{{pro\text{-}\acute{e}tale}}
\def\etale{{\acute{e}tale}}
\def\QCoh{\mathit{QCoh}}
\def\Ker{\mathop{\mathrm{Ker}}}
\def\Im{\mathop{\mathrm{Im}}}
\def\Coker{\mathop{\mathrm{Coker}}}
\def\Coim{\mathop{\mathrm{Coim}}}

% Boxtimes
%
\DeclareMathSymbol{\boxtimes}{\mathbin}{AMSa}{"02}

%
% Macros for moduli stacks/spaces
%
\def\QCohstack{\mathcal{QC}\!\mathit{oh}}
\def\Cohstack{\mathcal{C}\!\mathit{oh}}
\def\Spacesstack{\mathcal{S}\!\mathit{paces}}
\def\Quotfunctor{\mathrm{Quot}}
\def\Hilbfunctor{\mathrm{Hilb}}
\def\Curvesstack{\mathcal{C}\!\mathit{urves}}
\def\Polarizedstack{\mathcal{P}\!\mathit{olarized}}
\def\Complexesstack{\mathcal{C}\!\mathit{omplexes}}
% \Pic is the operator that assigns to X its picard group, usage \Pic(X)
% \Picardstack_{X/B} denotes the Picard stack of X over B
% \Picardfunctor_{X/B} denotes the Picard functor of X over B
\def\Pic{\mathop{\mathrm{Pic}}\nolimits}
\def\Picardstack{\mathcal{P}\!\mathit{ic}}
\def\Picardfunctor{\mathrm{Pic}}
\def\Deformationcategory{\mathcal{D}\!\mathit{ef}}
\usepackage{cleveref}
\crefname{originallemma}{original lemma}{original lemmas}
\Crefname{originallemma}{Original lemma}{Original lemmas}
\makeatletter
\newlabel{lemma-higher-direct-image}{{0720}{source}{}{sourcepointer-1}{}}
\newlabel{lemma-higher-direct-image@cref}{{[originallemma][0][]0720}{[1][0][]source}{}{}{}}
\makeatother
\crefname{originallemma}{original lemma}{original lemmas}
\Crefname{originallemma}{Original lemma}{Original lemmas}
\makeatletter
\newlabel{lemma-vanishing-quasi-separated}{{072C}{source}{}{sourcepointer-2}{}}
\newlabel{lemma-vanishing-quasi-separated@cref}{{[originallemma][0][]072C}{[1][0][]source}{}{}{}}
\makeatother
\crefname{originalproposition}{original proposition}{original propositions}
\Crefname{originalproposition}{Original proposition}{Original propositions}
\makeatletter
\newlabel{proposition-vanishing}{{072B}{source}{}{sourcepointer-3}{}}
\newlabel{proposition-vanishing@cref}{{[originalproposition][0][]072B}{[1][0][]source}{}{}{}}
\makeatother
\crefname{originallemma}{original lemma}{original lemmas}
\Crefname{originallemma}{Original lemma}{Original lemmas}
\makeatletter
\newlabel{sites-cohomology-lemma-apply-Leray}{{0733}{source}{}{sourcepointer-4}{}}
\newlabel{sites-cohomology-lemma-apply-Leray@cref}{{[originallemma][0][]0733}{[1][0][]source}{}{}{}}
\makeatother
\crefname{originallemma}{original lemma}{original lemmas}
\Crefname{originallemma}{Original lemma}{Original lemmas}
\makeatletter
\newlabel{spaces-properties-lemma-pushforward-etale-base-change-modules}{{03LX}{source}{}{sourcepointer-5}{}}
\newlabel{spaces-properties-lemma-pushforward-etale-base-change-modules@cref}{{[originallemma][0][]03LX}{[1][0][]source}{}{}{}}
\makeatother
\crefname{originallemma}{original lemma}{original lemmas}
\Crefname{originallemma}{Original lemma}{Original lemmas}
\makeatletter
\newlabel{spaces-properties-lemma-quasi-compact-affine-cover}{{03H6}{source}{}{sourcepointer-6}{}}
\newlabel{spaces-properties-lemma-quasi-compact-affine-cover@cref}{{[originallemma][0][]03H6}{[1][0][]source}{}{}{}}
\makeatother
\title{A uniform bound on the degrees in which the higher direct images vanish, stable under arbitrary base change}
\author{The Stacks Project Authors}
\date{stacks-project:spaces-cohomology.tex:local-snapshot-20261012; GNU FDL 1.2 or later}
\begin{document}
\maketitle
\noindent\textit{Editorial note.} All delivered mathematical text is the author's own, copied from the pinned Stacks Project chapter file spaces-cohomology.tex; no statement or proof step was rewritten, reordered or removed.\par\medskip
\noindent\textit{Editorial note.} Preamble provenance: the pinned chapter file opens with \textbackslash{}input\{preamble\} and carries no preamble text of its own, so the project's single Stacks-wide preamble staged at sources/stacks-spaces-cohomology/preamble.tex is excerpted verbatim as three spans, lines 8-12 (verbatim package and the project's own reference/slogan/history display environments), lines 14-18 (Xy-pic setup) and lines 154-217 (theorem declarations and mathematical macros). The wrapper class and the amsmath/amssymb/amsthm/graphicx base packages are supplied by the assembler. No macro was invented or redefined.\par\medskip
\noindent\textit{Editorial note.} Same-chapter pointers are located at their \textbackslash{}label lines inside the pinned chapter file; cross-chapter pointers are located in the project's own tag list staged in the snapshot (sources/stacks-chow/tags.txt), at the verbatim lines carrying their canonical tags.\par\medskip
\noindent\textit{Allowed standard prerequisites (editorial).}\begin{itemize}
\item An affine etale cover of a quasi-compact algebraic space, and base change of pushforwards along it. Properties of Spaces, labels spaces-properties-lemma-quasi-compact-affine-cover (invoked at proof line 1485) and spaces-properties-lemma-pushforward-etale-base-change-modules (invoked at proof line 1493). They reduce the statement to the case where the target \$Y\$ is an affine scheme, since the higher direct images restrict to those of the base changed morphism.
\item Quasi-coherence of the higher direct images along a quasi-compact quasi-separated morphism. Same chapter spaces-cohomology.tex, label lemma-higher-direct-image (invoked at proof line 1500). With an affine target \$Y'\$ it makes \$R\textasciicircum{}if'\_*\textbackslash{}mathcal\{F\}'\$ a quasi-coherent module whose global sections compute \$H\textasciicircum{}i(X', \textbackslash{}mathcal\{F\}')\$.
\item The Leray identification and vanishing on an affine space. Cohomology on Sites, label sites-cohomology-lemma-apply-Leray (invoked at proof line 1503) and Same chapter spaces-cohomology.tex, label proposition-vanishing (invoked at proof line 1506). Together they turn the vanishing of the higher direct images into the vanishing of \$H\textasciicircum{}i(X', \textbackslash{}mathcal\{F\}')\$ over an affine base.
\item The collection of covers and degrees attached to a quasi-separated space. Same chapter spaces-cohomology.tex, label lemma-vanishing-quasi-separated (invoked at proof lines 1510 and 1515). Its choices for \$X\$, base changed to \$X'\$, give the vanishing \$H\textasciicircum{}i(X', \textbackslash{}mathcal\{F\}') = 0\$ for \$i \textbackslash{}geq \textbackslash{}max(p + d\_p)\$, which is the bound asserted.
\end{itemize}\medskip
\ifdefined\hypertarget\hypertarget{sourcepointer-1}{}\fi
\noindent\textit{Source reference.} Original lemma 0720: sources/stacks-spaces-cohomology/pinned-original-source.tex, source lines 339--339. Navigation pointer for \textbackslash{}ref\{lemma-higher-direct-image\} used inside the counted span (line 1500); the argument that consumes it is retained in the span.\par\medskip
\ifdefined\hypertarget\hypertarget{sourcepointer-2}{}\fi
\noindent\textit{Source reference.} Original lemma 072C: sources/stacks-spaces-cohomology/pinned-original-source.tex, source lines 1385--1385. Navigation pointer for \textbackslash{}ref\{lemma-vanishing-quasi-separated\} used inside the counted span (line 1510); the argument that consumes it is retained in the span.\par\medskip
\ifdefined\hypertarget\hypertarget{sourcepointer-3}{}\fi
\noindent\textit{Source reference.} Original proposition 072B: sources/stacks-spaces-cohomology/pinned-original-source.tex, source lines 1333--1333. Navigation pointer for \textbackslash{}ref\{proposition-vanishing\} used inside the counted span (line 1506); the argument that consumes it is retained in the span.\par\medskip
\ifdefined\hypertarget\hypertarget{sourcepointer-4}{}\fi
\noindent\textit{Source reference.} Original lemma 0733: sources/stacks-chow/tags.txt, source lines 9047--9047. Navigation pointer for \textbackslash{}ref\{sites-cohomology-lemma-apply-Leray\} used inside the counted span (line 1503); the argument that consumes it is retained in the span.\par\medskip
\ifdefined\hypertarget\hypertarget{sourcepointer-5}{}\fi
\noindent\textit{Source reference.} Original lemma 03LX: sources/stacks-chow/tags.txt, source lines 4729--4729. Navigation pointer for \textbackslash{}ref\{spaces-properties-lemma-pushforward-etale-base-change-modules\} used inside the counted span (line 1493); the argument that consumes it is retained in the span.\par\medskip
\ifdefined\hypertarget\hypertarget{sourcepointer-6}{}\fi
\noindent\textit{Source reference.} Original lemma 03H6: sources/stacks-chow/tags.txt, source lines 4555--4555. Navigation pointer for \textbackslash{}ref\{spaces-properties-lemma-quasi-compact-affine-cover\} used inside the counted span (line 1485); the argument that consumes it is retained in the span.\par\medskip
\begin{lemma}
\label{lemma-vanishing-higher-direct-images}
Let $S$ be a scheme. Let $f : X \to Y$ be a
morphism of algebraic spaces over $S$.
Assume that
\begin{enumerate}
\item $f$ is quasi-compact and quasi-separated, and
\item $Y$ is quasi-compact.
\end{enumerate}
Then there exists an integer $n(X \to Y)$ such that
for any algebraic space $Y'$, any morphism $Y' \to Y$
and any quasi-coherent sheaf $\mathcal{F}'$ on $X' = Y' \times_Y X$
the higher direct images $R^if'_*\mathcal{F}'$ are zero for
$i \geq n(X \to Y)$.
\end{lemma}

\begin{proof}
Let $V \to Y$ be a surjective \'etale morphism where $V$ is an affine
scheme, see
Properties of Spaces, Lemma
\ref{spaces-properties-lemma-quasi-compact-affine-cover}.
Suppose we prove the result for the base change $f_V : V \times_Y X \to V$.
Then the result holds for $f$ with $n(X \to Y) = n(X_V \to V)$.
Namely, if $Y' \to Y$ and $\mathcal{F}'$ are as in the lemma, then
$R^if'_*\mathcal{F}'|_{V \times_Y Y'}$ is equal to
$R^if'_{V, *}\mathcal{F}'|_{X'_V}$ where
$f'_V : X'_V = V \times_Y Y' \times_Y X \to V \times_Y Y' = Y'_V$, see
Properties of Spaces,
Lemma \ref{spaces-properties-lemma-pushforward-etale-base-change-modules}.
Thus we may assume that $Y$ is an affine scheme.

\medskip\noindent
Moreover, to prove the vanishing for all $Y' \to Y$ and
$\mathcal{F}'$ it suffices to do so when $Y'$ is an affine scheme.
In this case, $R^if'_*\mathcal{F}'$ is quasi-coherent by
Lemma \ref{lemma-higher-direct-image}.
Hence it suffices to prove that $H^i(X', \mathcal{F}') = 0$, because
$H^i(X', \mathcal{F}') = H^0(Y', R^if'_*\mathcal{F}')$ by
Cohomology on Sites, Lemma \ref{sites-cohomology-lemma-apply-Leray}
and the vanishing of higher cohomology of quasi-coherent sheaves
on affine algebraic spaces
(Proposition \ref{proposition-vanishing}).

\medskip\noindent
Choose $U \to X$, $d$, $V_p \to U_p$ and $d_p$ as in
Lemma \ref{lemma-vanishing-quasi-separated}.
For any affine scheme $Y'$ and morphism $Y' \to Y$ denote
$X' = Y' \times_Y X$, $U' = Y' \times_Y U$, $V'_p = Y' \times_Y V_p$.
Then $U' \to X'$, $d' = d$, $V'_p \to U'_p$ and $d'_p = d$
is a collection of choices as in
Lemma \ref{lemma-vanishing-quasi-separated}
for the algebraic space $X'$ (details omitted).
Hence we see that $H^i(X', \mathcal{F}') = 0$ for $i \geq \max(p + d_p)$
and we win.
\end{proof}

\end{document}
